%%
%% Copyright (C) 2017 by Ivan Gankevich <igankevich@ya.ru>
%%
%% This file may be distributed and/or modified under the conditions of
%% the LaTeX Project Public License, either version 1.3c of this license
%% or (at your option) any later version.  The latest version of this
%% license is in:
%% 
%%    http://www.latex-project.org/lppl.txt
%% 
%% and version 1.3c or later is part of all distributions of LaTeX version
%% 2005/12/01 or later.
%%

\documentclass{ltxdoc}

\usepackage{xcolor}
\definecolor{spbuTerracotta}{cmyk}{.08,.91,.92,.33}
\definecolor{spbuGray}{cmyk}{.21,.11,.09,.22}
\definecolor{spbuDarkGray2}{HTML}{5F7177}
\definecolor{spbuWhite1}{RGB}{245,246,245}
\usepackage[colorlinks=true,
linkcolor=spbuTerracotta,
menucolor=spbuTerracotta,
pagecolor=spbuTerracotta,
urlcolor=spbuTerracotta]{hyperref}
\usepackage{listings}
\lstset{%
	language=[LaTeX]{TeX},
	basicstyle=\ttfamily,
	keywordstyle=\color{spbuTerracotta},
	commentstyle=\color{spbuGray},
	stringstyle=\color{spbuDarkGray2},
	backgroundcolor=\color{spbuWhite1},
	numbers=none,
	showspaces=false,
	showstringspaces=false,
	showtabs=false,
	frame=none,
	breaklines=true,
	breakatwhitespace=false,
	morekeywords={usetheme,institute,maketitle,setdefaultlanguage%
		,titlegraphic,othergraphic,leftcolumnwidth,rightcolumnwidth,%
		middlecolumnwidth,includegraphics},
}
\usepackage{metalogo}
\usepackage{fontspec}
\setmainfont[
	Extension=.otf,
	UprightFont=*-Regular,
	BoldFont=*-Bold,
	ItalicFont=*-Italic,
	Mapping=tex-text
]{OldStandard}
\setromanfont[
	Extension=.otf,
	UprightFont=*-Regular,
	BoldFont=*-Bold,
	ItalicFont=*-Italic,
	Mapping=tex-text
]{OldStandard}
\setsansfont[
	Extension=.ttf,
	UprightFont=*-Regular,
	BoldFont=*-Bold,
	ItalicFont=*-Italic,
	BoldItalicFont=*-BoldItalic,
	Mapping=tex-text
]{OpenSans}
\setmonofont[
	Extension=.otf,
	UprightFont=*-Regular,
	BoldFont=*-Bold
]{FiraMono}

\title{Saint Petersburg: \LaTeX{} Beamer theme for SPbU}
\author{Ivan Gankevich}
\date{20 Nov 2017}

\begin{document}

\maketitle
\tableofcontents

\section{Introduction}

A small theme that incorporates university colours and fonts from
\href{http://pr.spbu.ru/}{the official web-site}.

\subsection{Installation}

On Linux/MacOS clone \href{https://github.com/igankevich/SaintPetersburg}{theme
repository} and type \verb=make install= to install everything to standard TeX
Live locations. Alternatively, just copy all \verb=*.sty= files into your
project directory so that \LaTeX{} can find them.  In order to show
university's logo in the background of the title slide, you need to download
its medium size version from
\href{http://pr.spbu.ru/images/simvolika/logo/CoA_Medium.eps}{the official
web-site}. 

\subsection{Usage}

Saint Petersburg theme can be compiled by \LaTeX{} or \XeLaTeX{}. Here is the
minimal working example:
\begin{lstlisting}
\documentclass[aspectratio=169]{beamer}
% add \usepackage{beamerposter} for the poster

% XeTeX
\usepackage{polyglossia}
\setdefaultlanguage{english}
% or \setdefaultlanguage{russian}

% LaTeX
% \usepackage[utf8]{inputenc}
% \usepackage[english]{babel}
% or \usepackage[english,russian]{babel}

\usetheme{SaintPetersburg}

\title{Saint Petersburg \LaTeX{} Beamer theme}
\author{Ivan Gankevich}
\institute{Saint Petersburg State University}
\date{20 Nov 2017}

\begin{document}
\frame{\titlepage}
\end{document}
\end{lstlisting}

\noindent An example of poster.
\begin{lstlisting}
\documentclass[aspectratio=169]{beamer}
\usepackage{beamerposter} for the poster

% XeTeX
\usepackage{polyglossia}
\setdefaultlanguage{english}
% or \setdefaultlanguage{russian}

% LaTeX
% \usepackage[utf8]{inputenc}
% \usepackage[english]{babel}
% or \usepackage[english,russian]{babel}

\usetheme[poster]{SaintPetersburg}

\title{Saint Petersburg \LaTeX{} Beamer theme}
\author{Ivan Gankevich}
\institute{Saint Petersburg State University}
\date{20 Nov 2017}

\begin{document}
\frame{\titlepage}
\end{document}
\end{lstlisting}

\subsection{Usage}

\subsubsection{Theme options}

\DescribeMacro{nologo} Disables display of the university coat of arms in the
left top corner of the title page.\vspace{\baselineskip}

\noindent\DescribeMacro{poster} Adapts the theme for poster presentation:
enables section commands, increases font size and changes default font family
to serif.

\subsubsection{Theme macros}

\DescribeMacro{\titlegraphic}%
\DescribeMacro{\othergraphic}%
Defines a code to include in the left (right) top corner of the poster. The
following code places SPbU logo in both left and right corners.
\begin{lstlisting}
\titlegraphic{%
	\includegraphics[width=.7\linewidth]{spbu-CoA}}
\othergraphic{%
	\includegraphics[width=.7\linewidth]{spbu-CoA}}
\end{lstlisting}
Here the width equals to 70\% of the left (right) column width. Any code can be
specified instead of just including graphics.

\DescribeMacro{\leftcolumnwidth}%
\DescribeMacro{\rightcolumnwidth}%
\DescribeMacro{\middlecolumnwidth}%
Sets left, right, middle column widths repsectively. The right and left columns
contain logos and the middle column contains the title, authors and institute
name. In the following example, the widths are set to 20\%, 60\% and 20\%
fractions of the total line width.
\begin{lstlisting}
\leftcolumnwidth{.2\linewidth}
\middlecolumnwidth{.6\linewidth}
\rightcolumnwidth{.2\linewidth}
\end{lstlisting}

\subsection{License}

Saint Petersburg theme is licensed under the terms of the
\href{https://ctan.org/license/lppl1.3c}{\LaTeX{} Project Public License
version 1.3c}.

\section{Implementation}

\DocInput{beamercolorthemeSaintPetersburg.dtx}
\DocInput{beamerfontthemeSaintPetersburg.dtx}
\DocInput{beamerthemeSaintPetersburg.dtx}

\end{document}
